%=============================================================================
% APUNTES DE DERECHO REALES
% Derecho de Superficie, Usufructo y Derechos Reales de Disfrute
% Cátedra de Derechos Reales — Facultad de Derecho — UBA
% Compilación: 19 de septiembre de 2026
%=============================================================================

\documentclass[12pt,oneside]{book}

%--- CODIFICACIÓN ---
\usepackage[utf8]{inputenc}
\usepackage[T1]{fontenc}
\usepackage[spanish,es-nodecimaldot]{babel}

%--- GEOMETRÍA DE PÁGINA ---
\usepackage[
    top=2.5cm,
    bottom=2.5cm,
    left=3cm,
    right=2.5cm,
    headheight=15pt
]{geometry}

%--- TIPOGRAFÍA ---
\usepackage{ragged2e}

%--- MATEMÁTICA Y SÍMBOLOS ---
\usepackage{amsmath}
\usepackage{amssymb}
\usepackage{mathtools}

%--- IMÁGENES ---
\usepackage{graphicx}
\usepackage{float}
\usepackage{caption}
\usepackage{subcaption}

%--- TABLAS ---
\usepackage{booktabs}
\usepackage{array}
\usepackage{longtable}
\usepackage{tabularx}

%--- LISTAS ---
\usepackage{enumitem}

%--- COLORES Y CAJAS ---
\usepackage{xcolor}
\usepackage[most]{tcolorbox}

%--- ENCABEZADOS Y PIES ---
\usepackage{fancyhdr}

%--- HIPERENLACES Y METADATOS ---
\usepackage[
    colorlinks=true,
    linkcolor=blue!50!black,
    citecolor=green!40!black,
    urlcolor=blue!60!black,
    pdftitle={Apuntes de Derecho Reales},
    pdfauthor={Isaac Edgar Camacho Ocampo},
    pdfsubject={Apuntes universitarios},
    bookmarks=true
]{hyperref}

%=============================================================================
% CONFIGURACIÓN GENERAL
%=============================================================================

\setlength{\parindent}{0pt}
\setlength{\parskip}{0.5em}

\setlist[itemize]{
    topsep=0.3em,
    itemsep=0.2em
}

\setlist[enumerate]{
    topsep=0.3em,
    itemsep=0.2em
}

\clubpenalty=10000
\widowpenalty=10000

\setcounter{tocdepth}{2}

%=============================================================================
% COMANDOS MATEMÁTICOS PERSONALIZADOS
%=============================================================================

\newcommand{\abs}[1]{\left|#1\right|}
\newcommand{\norm}[1]{\left\lVert#1\right\rVert}
\newcommand{\vect}[1]{\mathbf{#1}}
\newcommand{\deriv}[2]{\frac{d#1}{d#2}}
\newcommand{\pderiv}[2]{\frac{\partial#1}{\partial#2}}

%=============================================================================
% DEFINICIÓN DE CAJAS ACADÉMICAS
%=============================================================================

\newtcolorbox{definition}[1]{
    enhanced,
    breakable,
    title={Definición: #1},
    fonttitle=\bfseries,
    colback=blue!3,
    colframe=blue!50!black,
    boxrule=0.8pt,
    arc=2mm
}

\newtcolorbox{important}{
    enhanced,
    breakable,
    title={Importante},
    fonttitle=\bfseries,
    colback=yellow!8,
    colframe=orange!70!black,
    boxrule=0.8pt,
    arc=2mm
}

\newtcolorbox{example}[1]{
    enhanced,
    breakable,
    title={Ejemplo: #1},
    fonttitle=\bfseries,
    colback=green!4,
    colframe=green!40!black,
    boxrule=0.8pt,
    arc=2mm
}

\newtcolorbox{formula}[1]{
    enhanced,
    breakable,
    title={Fórmula/Regla: #1},
    fonttitle=\bfseries,
    colback=purple!4,
    colframe=purple!50!black,
    boxrule=0.8pt,
    arc=2mm
}

\newtcolorbox{exercise}[1]{
    enhanced,
    breakable,
    title={Ejercicio: #1},
    fonttitle=\bfseries,
    colback=gray!5,
    colframe=gray!60!black,
    boxrule=0.8pt,
    arc=2mm
}

\newtcolorbox{note}{
    enhanced,
    breakable,
    title={Nota},
    fonttitle=\bfseries,
    colback=white,
    colframe=gray!50,
    boxrule=0.6pt,
    arc=2mm
}

%=============================================================================
% ENCABEZADOS Y PIES DE PÁGINA
%=============================================================================

\pagestyle{fancy}
\fancyhf{}

\fancyhead[L]{\nouppercase{\leftmark}}
\fancyhead[R]{Derechos Reales}

\fancyfoot[C]{\thepage}

\renewcommand{\headrulewidth}{0.4pt}
\renewcommand{\footrulewidth}{0pt}

\fancypagestyle{plain}{
    \fancyhf{}
    \fancyfoot[C]{\thepage}
    \renewcommand{\headrulewidth}{0pt}
}

%=============================================================================
% INICIO DEL DOCUMENTO
%=============================================================================

\begin{document}

%=============================================================================
% PORTADA
%=============================================================================

\begin{titlepage}

\thispagestyle{empty}

\begin{center}

\vspace*{1cm}

{\large \textsc{Universidad de Buenos Aires}}

\vspace{0.2cm}

{\large \textsc{Facultad de Derecho}}

\vspace{0.2cm}

{\large Cátedra de Derechos Reales}

\vspace{3cm}

{\Huge \textbf{Apuntes de Derecho Reales}}

\vspace{0.8cm}

{\LARGE Derecho de Superficie, Usufructo y Derechos Reales de Disfrute}

\vspace{0.8cm}

\textit{El estudio riguroso de los derechos reales construye las herramientas conceptuales para resolver conflictos patrimoniales reales.}

\vspace{1.2cm}

\rule{0.70\textwidth}{0.5pt}

\vspace{0.5cm}

{\large Apuntes de estudio}

\vspace{0.5cm}

\rule{0.70\textwidth}{0.5pt}

\vfill

{\large Isaac Edgar Camacho Ocampo}

\vspace{0.3cm}

{\large 2026}

\end{center}

\vfill

\begin{flushleft}
\textit{Este es un modesto aporte a los estudiantes de ``Derecho Reales'' no pretende sustituir el material oficial de la materia.}
\end{flushleft}

\end{titlepage}

%=============================================================================
% ÍNDICE DE CONTENIDOS
%=============================================================================

\tableofcontents

%=============================================================================
% INTRODUCCIÓN
%=============================================================================

\chapter*{Introducción}

\section*{Sobre esta compilación}

Esta compilación recoge los contenidos desarrollados en la clase sobre derechos reales que se desprenden del dominio. El material ha sido estructurado siguiendo un hilo conductor único: la manera en que una familia accede a la vivienda, la financia, la transmite o la preserva para generaciones futuras, empleando para ello las diversas figuras que el derecho real pone a disposición.

\section*{Bloques temáticos}

Los siete bloques que conforman estos apuntes abarcan:

\begin{enumerate}
\item El repaso del dominio y los objetivos del curso
\item El derecho de superficie en profundidad, como herramienta de financiamiento inmobiliario
\item El usufructo, la reversión y las herramientas de planificación sucesoria
\item Las cláusulas societarias trasladadas al condominio
\item El condominio, la sucesión y las sociedades
\item El derecho de uso y el derecho de habitación
\item Un trabajo práctico integrador
\end{enumerate}

\section*{Orientación pedagógica}

Cada figura jurídica se presenta tanto en sus aspectos teóricos como en sus aplicaciones prácticas. El objetivo no es ofrecerle a quien estudia un resumen exhaustivo, sino proporcionarle las herramientas conceptuales necesarias para comprender, diseñar e instrumentar soluciones reales frente a problemas concretos de vivienda, financiamiento y planificación patrimonial.

%=============================================================================
% BLOQUE 1 — REPASO Y CONTEXTO DEL CURSO
%=============================================================================

\chapter{Repaso y contexto del curso}

\section{Repaso del dominio y objetivos del curso}

El estudio de los derechos reales que se desprenden del dominio no es casual en la estructuración de esta materia. El método adoptado responde a una lógica pedagógica específica: se estudia primero el dominio pleno para que, a partir de él, puedan comprenderse los derechos reales sobre cosa ajena. Esta progresión permite también proyectar las nuevas figuras vinculadas a la tecnología ---como el blockchain--- que, en paralelo, se abordan en el curso.

\begin{important}
El dominio no se limita a los bienes inmuebles. La misma lógica se aplica al derecho societario, en particular a la posibilidad de constituir usufructo sobre cuotas sociales o sobre acciones, algo que la ley de sociedades vigente sí permite en Argentina, aunque todavía existen algunas figuras que el derecho argentino no admite en ese terreno.
\end{important}

La estructura del curso responde, en simultáneo, a dos frentes de trabajo: lo rígido ---la estructura clásica y tradicional de los derechos reales tal como está regulada--- y lo tecnológico ---las nuevas posibilidades que ofrece la tecnología para los mismos institutos. Como todavía no se sabe con certeza hacia dónde evolucionará la implementación tecnológica, es necesario enseñar ambos caminos: el analógico, porque los estudiantes van a tener que trabajar con él en la práctica profesional actual, y el tecnológico, porque representa el futuro de la disciplina.

%=============================================================================
% BLOQUE 2 — EL DERECHO DE SUPERFICIE EN PROFUNDIDAD
%=============================================================================

\chapter{El derecho de superficie en profundidad}

\section{Estructura del derecho de superficie}

\begin{definition}{Derecho de Superficie}
Derecho real que permite construir, forestar o plantar sobre suelo, subsuelo o espacio aéreo ajeno, comprendiendo también lo construido, lo plantado, lo edificado o lo forestado sobre esas porciones. El titular de este derecho ---el superficiario--- puede enajenarlo a terceros, transmitiendo su posición como derecho superficiario.
\end{definition}

El derecho de superficie se extingue cuando se extingue el derecho de superficie del cual depende. Un aspecto central ---que debe quedar con claridad--- es que quien constituye un derecho de superficie sobre un inmueble adquiere, sobre lo construido, la posibilidad de enajenar ese derecho a un tercero.

\section{Financiamiento de vivienda mediante el derecho de superficie}

El derecho de superficie constituye una herramienta especialmente potente para el desarrollo y financiamiento inmobiliario. A modo de ilustración, considérese el siguiente esquema:

\begin{example}{Financiamiento de vivienda mediante derecho de superficie}

Una persona ---la propietaria del terreno--- constituye a favor de otra ---el superficiario--- un derecho de superficie sobre su terreno, por un plazo de 30 a 50 años. El superficiario construye sobre ese terreno veinte unidades habitacionales.

Sobre esas veinte unidades, el superficiario puede vender a terceros el derecho de superficie en forma financiada: los compradores entran con una entrega inicial y pagan cuotas mensuales durante, por ejemplo, diez años.

Con lo producido por la venta financiada de las veinte unidades, el superficiario recupera lo invertido en la construcción mucho antes de que se cumpla el plazo total de su propio derecho de superficie (30 a 50 años). De esta manera, genera un derecho a la vivienda alternativo: permite a compradores que no pueden acceder a la titularidad de dominio pleno adquirir, en cambio, un derecho de superficie financiado, con una seguridad jurídica muy superior a la de un simple alquiler.

\end{example}

\section{La hipoteca sobre el derecho de superficie}

Cuando se otorga una hipoteca para garantizar la venta financiada de una unidad superficiaria, la garantía se constituye sobre el propio derecho de superficie del comprador, nunca sobre el dominio del terreno, que sigue perteneciendo a su dueña original.

\begin{important}
Quien se asegura mediante hipoteca es el dueño del derecho de superficie, no la propietaria del terreno, quien no participa de esa garantía. En el ejemplo de las veinte unidades, cada uno de los compradores adquiere una de ellas sobre la base de un derecho de superficie, y esa adquisición se financia con una hipoteca constituida sobre ese mismo derecho de superficie ---nunca sobre el dominio del terreno, que permanece en cabeza de su titular original.
\end{important}

El superficiario originario puede incorporar dentro de ese derecho de superficie subcontratado múltiples cláusulas contractuales para proteger su posición. Por ejemplo:

\begin{itemize}
\item Que el comprador no pueda alquilar la unidad mientras no cancele la hipoteca
\item Que no pueda disponer de ella hasta el pago total del precio
\item Que no pueda vender a personas que el vendedor no acepte
\end{itemize}

Todas esas condiciones quedan incorporadas y son oponibles dentro del propio derecho de superficie.

\section{Subcontratación, plazos y reversión del derecho}

Un punto frecuente de confusión es qué ocurre cuando el derecho de superficie originario (por ejemplo, de 50 años) se subcontrata por un plazo menor (por ejemplo, 20 años) con los compradores finales.

\begin{formula}{Regla de extinción y reversión}
Cuando el derecho de superficie subcontratado se extingue (después de 20 años), no revierte a la propietaria original del terreno, sino al superficiario intermedio que lo había constituido ---quien conserva su propio derecho de superficie por el plazo restante. El principio general que rige toda la figura es que \textbf{nadie puede adquirir más derechos de los que quien transmite estaba en condiciones de transmitir}.
\end{formula}

Los adquirentes de las unidades superficiarias están en una posición distinta a la de un simple locatario: aunque no sean titulares de dominio pleno, durante el plazo de su derecho pueden vender su derecho, alquilarlo o hipotecarlo, es decir, disponer de él con una amplitud que un simple locatario nunca tiene.

Extinguido finalmente el derecho de superficie originario en su totalidad (cumplidos los 30 a 50 años, salvo pacto en contrario), la propiedad vuelve a su titular original. Sin embargo, en el propio contrato de constitución pueden pactarse mil situaciones alternativas para ese momento:

\begin{itemize}
\item Que el superficiario pague una suma de dinero a cambio de que se le transmita el dominio pleno
\item Que se prorrogue el derecho de superficie por un plazo adicional
\item Que se reconozca al superficiario una preferencia u opción de compra sobre el inmueble
\item Que ambas partes queden en condominio sobre todo o parte del bien
\end{itemize}

\section{Superficie versus usufructo: la extinción}

Una diferencia central entre el derecho de superficie y el usufructo radica en sus causales de extinción:

\begin{important}
Cuando una persona se reserva el usufructo de un bien y cede la nuda propiedad (por ejemplo, a un hijo), ese usufructo se extingue con la muerte del usufructuario. El derecho de superficie, en cambio, no se extingue por fallecimiento del superficiario: se extingue por el vencimiento del plazo pactado o por las demás causales contractuales aplicables, pero nunca por la sola muerte de su titular.
\end{important}

%=============================================================================
% BLOQUE 3 — USUFRUCTO, REVERSIÓN Y PLANIFICACIÓN SUCESORIA
%=============================================================================

\chapter{Usufructo, reversión y planificación sucesoria}

\section{El derecho de reversión en las donaciones}

\begin{definition}{Derecho de Reversión}
Cláusula que puede incorporarse a una donación y que permite al donante recuperar el bien donado si el donatario fallece antes que él, siempre que ese fallecimiento ocurra dentro del plazo de diez años contados desde la donación. Vencido ese plazo, la reversión deja de ser exigible.
\end{definition}

\begin{important}
El derecho de reversión solo puede ejercerse si el donatario fallece dentro de un plazo de diez años desde la donación. Transcurrido ese plazo, el derecho de reversión ya no es una opción disponible, aunque se hubiera pactado expresamente.
\end{important}

Una advertencia práctica importante: es riesgoso pactar solamente el derecho de reversión sin reservarse además el usufructo del bien. Si el donante fallece antes de que se cumplan esos diez años, pierde toda posibilidad de usar la cosa durante ese lapso, ya que se desprendió tanto de la nuda propiedad como del uso.

\section{Riesgos de la donación como planificación sucesoria}

La práctica frecuente de donar bienes en vida para evitar la sucesión merece un análisis crítico. La donación puede ser una herramienta razonable cuando la persona está en una edad avanzada o en una situación de salud que hace previsible el desenlace, pero no como recurso preventivo general y anticipado.

Entre los riesgos concretos que genera la donación se cuentan:

\begin{itemize}
\item \textbf{Exposición a contingencias del donatario}: una vez que el bien ingresa al patrimonio del hijo donatario, queda expuesto a las contingencias propias de ese patrimonio. Por ejemplo, si el hijo atraviesa un divorcio, el bien puede verse afectado por una medida de no innovar sobre bienes propios, aun cuando el donante hubiera conservado el usufructo.

\item \textbf{Limitación del acceso al crédito}: el hijo, al ser titular registral de un inmueble (aunque no lo use, porque el usufructo sigue en cabeza del donante), puede ver limitado su acceso al crédito bancario, porque a los ojos de la entidad financiera ya es propietario de una vivienda.

\item \textbf{Falta de flexibilidad}: la donación es un acto difícil de revocar y carece de flexibilidad frente a cambios en la situación familiar o económica.
\end{itemize}

Antes de recurrir automáticamente a la donación, al usufructo o a otras figuras "tradicionales", conviene evaluar si existen otras herramientas mejor adaptadas a la situación concreta de cada familia.

\section{Cesión del usufructo y renta vitalicia}

El usufructuario puede ceder su derecho a un tercero. Sin embargo, quien adquiere el usufructo cedido queda sujeto al mismo límite temporal que el usufructuario original, de modo que el usufructo cedido se extingue con el fallecimiento del usufructuario originario, no con el del cesionario.

\begin{definition}{Renta Vitalicia}
Contrato aleatorio cuyo resultado económico depende de un hecho incierto: la duración de la vida del beneficiario. La renta vitalicia implica una transacción donde una parte transmite la propiedad de un bien a cambio de recibir una renta periódica durante toda su vida.
\end{definition}

\begin{important}
La renta vitalicia constituye un contrato aleatorio que puede ser excelente o pésimo negocio según variables que no pueden conocerse de antemano. Nadie puede saber con certeza cuánto vivirá el rentista. Además, quien recibe la renta puede ver comprometida su capacidad de pago si las circunstancias externas cambian (por ejemplo, si aumentan los costos de un geriátrico que se comprometió a pagar).
\end{important}

Esta reflexión se aplica a cualquier decisión patrimonial sujeta a variables externas. La volatilidad de los mercados, la incertidumbre sobre el futuro y la imposibilidad de controlar ciertos factores exigen una actitud de cautela y revisión constante de las decisiones tomadas.

%=============================================================================
% BLOQUE 4 — CLÁUSULAS SOCIETARIAS APLICADAS A DERECHOS REALES
%=============================================================================

\chapter{Cláusulas societarias aplicadas a derechos reales}

\section{Cláusulas de arrastre (Drag-Along) y de acompañamiento (Tag-Along)}

Un problema típico del derecho societario es la situación de los socios minoritarios frente a las decisiones de los mayoritarios. Para resolver esta cuestión, existen dos cláusulas incorporadas expresamente al ordenamiento a partir de la Ley de Sociedades por Acciones Simplificadas (Ley de SAS):

\begin{definition}{Derecho de Arrastre (Drag-Along)}
Cláusula que obliga a los socios minoritarios a vender también su participación cuando el socio principal decide vender la suya, de modo que el comprador pueda adquirir el 100\% del capital y los minoritarios no queden atrapados en una sociedad cuyo control cambió de manos sin ellos.
\end{definition}

\begin{definition}{Derecho de Acompañamiento (Tag-Along)}
Cláusula que otorga a los minoritarios la posibilidad ---no la obligación--- de vender su participación en las mismas condiciones y al mismo valor que se le pagó al socio principal, si así lo desean. Si no quieren vender, no están obligados a hacerlo.
\end{definition}

Estas cláusulas resultan especialmente convenientes en el diseño de una startup: quien tuvo la idea original y conserva, por ejemplo, un 30\% o un 40\% del capital, puede verse en la práctica sin ningún poder de decisión real si el socio mayoritario decide vender y él no cuenta con una cláusula de acompañamiento que le permita salir en las mismas condiciones.

\section{Estas cláusulas aplicadas al condominio}

La lógica del arrastre y del acompañamiento puede trasladarse al condominio de manera creativa. En un condominio compuesto por tres personas con el 20\%, el 40\% y el 40\% del bien, basta con que una sola de ellas no quiera vender para que nadie esté dispuesto a comprar una propiedad con apenas el 99\% de sus partes libres.

Por eso puede proponerse que los condóminos firmen un pacto que replique la lógica del arrastre y del acompañamiento: si existe un comprador interesado en la totalidad del inmueble, todos los condóminos quedan obligados a vender; y si la mayoría de los condóminos decide vender, la minoría restante puede sumarse a esa venta en las mismas condiciones.

\begin{important}
Los pactos entre condóminos son oponibles entre las partes que los suscriben, pero en principio no resultan oponibles frente a terceros. Sin embargo, la ley admite que un pacto sea oponible a un tercero que haya tomado conocimiento efectivo de su existencia, de modo que basta con exhibirle el pacto a ese tercero para que, a su criterio, quede alcanzado por él.
\end{important}

Para instrumentar en la práctica la obligación de vender asumida en el pacto, suele recurrirse a un poder especial e irrevocable, otorgado al momento de firmar el pacto, que permite ejecutar la venta llegado el caso sin depender de la voluntad posterior del condómino remiso. Por eso los pactos resultan interesantes: sin ese poder especial irrevocable el condómino podría simplemente negarse a firmar cuando llegara el momento.

%=============================================================================
% BLOQUE 5 — CONDOMINIO, SUCESIÓN Y SOCIEDADES
%=============================================================================

\chapter{Condominio, sucesión y sociedades}

\section{División de condominio y comunidad hereditaria}

Existen dos situaciones que, aunque producen efectos prácticos similares, tienen una naturaleza jurídica distinta: el condominio propiamente dicho, nacido de un derecho real, y la comunidad hereditaria, que si bien a los efectos prácticos funciona como un condominio, no lo es desde el punto de vista técnico.

\begin{example}{Herencia indivisa y venta}

Considérese el caso de dieciséis herederos, de los cuales quince están de acuerdo en vender y uno se niega. En ese supuesto, la vía es la división judicial de la herencia: no existe otra forma de forzar la partición si no se logra un acuerdo extrajudicial, de modo que la cuestión queda directamente judicializada.

Cuando en cambio se trata de un condominio de derecho real y no de una comunidad hereditaria, la vía correspondiente es la acción de división de condominio, también por la vía judicial si no se logra un acuerdo entre las partes.

\end{example}

\section{Riesgos de comprar derechos no consolidados}

Un riesgo frecuente en el mercado inmobiliario actual es la compra de emprendimientos que todavía no cuentan con planos aprobados ni con un derecho de superficie efectivamente constituido, sino apenas con la expectativa de que, en el futuro, se les adjudique una unidad funcional.

\begin{important}
En esos casos suele ocurrir que, al momento de aprobarse finalmente los planos, sea necesario modificar el proyecto original o ceder una porción del terreno, lo que puede afectar directamente a quienes ya compraron con base en un proyecto que todavía no era jurídicamente firme.
\end{important}

\begin{example}{El departamento en desarrollo sin derecho constituido}

Quien está en la posición de comprador debe preguntarse: ¿A quién le estoy comprando? ¿Qué es lo que estoy comprando exactamente? Muchas veces lo que en realidad se adquiere no es un derecho real consolidado, sino apenas la cesión de un boleto o de una expectativa contractual.

Este riesgo no depende únicamente del nivel de sofisticación del comprador. Muchas personas que pagan sumas millonarias por este tipo de operaciones cuentan con asesores y escribanos, y aun así terminan comprando derechos que no están firmes, porque el proyecto no tiene los planos aprobados o porque, para lograr esa aprobación, será necesario modificar la construcción original.

\end{example}

\section{Incorporación de herederos y quórum en asambleas}

La relación entre condominio, comunidad hereditaria y sociedades se torna compleja cuando fallece un socio o accionista y sus herederos todavía no fueron declarados como tales. Mientras no exista declaratoria de herederos ni partición, no hay certeza definitiva sobre quién representa al causante dentro de la sociedad, y que la propia partición podría, incluso, adjudicar esa participación societaria a un heredero distinto del que hasta ese momento venía ejerciendo los derechos del fallecido.

\begin{example}{Exclusión de herederos en asamblea}

Una situación práctica que genera conflictos: la familia de un fallecido socio asiste, al día siguiente de velar al familiar, a una asamblea de la que los herederos son excluidos con el argumento de que todavía no existe declaratoria de herederos que los legitime para participar. Ese tipo de situaciones ---donde hasta ayer eras el tío y llorabas conmigo, y hoy en la asamblea te desconoce--- puede derivar en la adopción de decisiones sociales que perjudican los intereses del causante y de sus sucesores.

\end{example}

Respecto a las asambleas extraordinarias, es necesario distinguir entre quórum y mayoría:

\begin{definition}{Quórum en Asambleas}
Porcentaje de capital con derecho a voto que debe estar presente para que la asamblea pueda sesionar válidamente. En primera convocatoria se exige un quórum del 60\% de las acciones con derecho a voto; en segunda convocatoria ese quórum se reduce a un porcentaje sensiblemente menor.
\end{definition}

\begin{definition}{Mayoría en Asambleas}
Porcentaje de votos necesario, sobre los presentes, para aprobar una decisión.
\end{definition}

\begin{important}
El estatuto puede elevar los porcentajes mínimos legales de quórum, pero nunca hasta el punto de exigir una unanimidad encubierta.
\end{important}

Existe en la doctrina debate sobre el momento en que se considera incorporado un heredero a la sociedad. Hay tres posturas:

\begin{enumerate}
\item Al momento del fallecimiento del causante socio o accionista
\item Al momento de la declaratoria de herederos y la designación de administrador de la sucesión
\item Al momento de la partición definitiva de la herencia
\end{enumerate}

Se trata de una cuestión difícil de resolver en la práctica, además de ser un tema que está siendo objeto de tratamiento en el proyecto de reforma de la ley de sociedades.

%=============================================================================
% BLOQUE 6 — DERECHO DE USO Y DERECHO DE HABITACIÓN
%=============================================================================

\chapter{Derecho de uso y derecho de habitación}

\section{Diferencias entre el uso y la habitación}

Aunque parecen figuras poco utilizadas en la práctica, bien empleadas resultan sumamente útiles. Es importante identificar claramente las diferencias entre ambos derechos:

\begin{table}[H]
\centering
\small
\begin{tabularx}{\textwidth}{
    >{\raggedright\arraybackslash}p{0.22\textwidth}
    >{\raggedright\arraybackslash}X
    >{\raggedright\arraybackslash}X
}
\toprule
\textbf{Criterio} & \textbf{Derecho de uso} & \textbf{Derecho de habitación} \\
\midrule

\textbf{Titularidad} & Puede constituirse a favor de personas humanas. & Solo puede constituirse a favor de personas humanas. \\

\textbf{Objeto} & Puede recaer sobre cosas muebles e inmuebles (por ejemplo, un tractor). & Solo puede recaer sobre un inmueble. \\

\textbf{Amplitud} & Derecho más amplio: además de usar, permite gozar de los frutos que la cosa produce. & Derecho más reducido, limitado a morar en el inmueble. \\

\textbf{Vigencia} & Puede ser vitalicio o pactarse por un plazo determinado. & Puede ser vitalicio o temporal. \\

\textbf{Figura personal equivalente} & Puede constituirse también como derecho personal (por ejemplo, mediante un contrato de comodato). & Solo se admite como derecho real; no como derecho personal equivalente. \\

\bottomrule
\end{tabularx}
\end{table}

\subsection{Solidez jurídica de los derechos reales}

Un derecho real de habitación, ya inscripto, se sostiene por sí mismo ---salvo que se configure alguna causal de extinción--- sin necesidad de recurrir a los tribunales para hacerlo valer. Un derecho meramente personal (como un comodato), en cambio, depende de la autonomía de la voluntad contractual y puede verse discutido por terceros, obligando a litigar para hacerlo respetar.

\subsection{Preferencia de habitación sobre usufructo}

En ciertos casos, suele preferirse el derecho de habitación por sobre el usufructo. Si a la persona que cuidó de un familiar se le constituye un usufructo, esta podría alquilar la vivienda a un tercero sin habitarla ella misma, desvirtuando la finalidad original de asegurarle un lugar donde vivir.

El derecho de habitación, en cambio, está más directamente orientado a garantizar la morada efectiva del beneficiario, sin que por ello se vea perjudicado ni tampoco perjudique a los demás herederos.

\begin{example}{Derecho de habitación como gesto de gratitud}

Una familia, en reconocimiento a la mujer que había cuidado a su madre durante años y que tenía prácticamente su misma edad, le constituyó un derecho real de habitación sobre la vivienda familiar hace aproximadamente diez años.

Con el tiempo, la madre falleció y, más adelante, la propia beneficiaria del derecho de habitación desarrolló Alzheimer, quedando imposibilitada de seguir viviendo sola en esa casa. Ante esa situación, la familia decidió no ejercer más el derecho de habitación en especie, sino alquilar la vivienda y destinar lo producido por ese alquiler a solventar los gastos de un geriátrico para la beneficiaria.

El derecho real de habitación no necesariamente tiene que traducirse en la ocupación efectiva del inmueble por parte del beneficiario: puede reconvertirse, según las circunstancias, en una fuente de sustento económico equivalente, sin perder su función original de gratitud y protección hacia quien prestó años de cuidado.

\end{example}

%=============================================================================
% BLOQUE 7 — TRABAJO PRÁCTICO Y CIERRE DE LA CLASE
%=============================================================================

\chapter{Trabajo práctico integrador}

\section{Consigna del trabajo práctico}

Se pide a los estudiantes que planteen situaciones concretas ---no hipótesis abstractas--- en las que aplicarían cada uno de los derechos reales vistos hasta el momento (uso, habitación, usufructo, superficie, entre otros), incorporando las cláusulas que consideren pertinentes para cada caso.

\subsection{Modelo de respuesta}

La respuesta esperada sigue un formato sencillo y aplicado, formulado en primera persona:

\begin{quote}
``Yo constituiría un derecho real de uso, por ejemplo, porque es conveniente para la explotación de tal cosa; a ella le daría un usufructo por tal motivo; a él le daría el derecho de habitación por tal situación; acá constituiría un derecho real de superficie sobre el subsuelo por tal circunstancia.''
\end{quote}

Se espera brevedad y aplicabilidad práctica, no extensión. El trabajo debe ser un ejercicio breve aplicado a casos concretos y propios de cada estudiante, no un trabajo extenso de veinte o cuarenta páginas.

\subsection{Objetivo del trabajo práctico}

El trabajo práctico no tiene como finalidad principal la acreditación de un puntaje, sino verificar, tanto para el estudiante como para la cátedra, si los contenidos vistos hasta el momento fueron efectivamente comprendidos y pueden aplicarse a situaciones reales y concretas.

%=============================================================================
% CUADRO CORNELL
%=============================================================================

\chapter*{Cuadro de síntesis Cornell}
\addcontentsline{toc}{chapter}{Cuadro de síntesis Cornell}

\thispagestyle{plain}

\vspace{1cm}

\section*{Método de estudio: Preguntas y respuestas clave}

\begin{table}[H]
\centering
\renewcommand{\arraystretch}{1.8}
\begin{tabularx}{\textwidth}{
    >{\raggedright\arraybackslash}p{0.28\textwidth}
    >{\raggedright\arraybackslash}X
}
\toprule
\textbf{Preguntas / Palabras clave} &
\textbf{Notas / Respuestas / Explicaciones} \\
\midrule

\textbf{¿Qué es el derecho de superficie y sobre qué puede recaer?} &
Derecho real que permite construir, forestar o plantar sobre suelo, subsuelo o espacio aéreo ajeno, comprendiendo lo construido, plantado, edificado o forestado sobre esas porciones. El superficiario puede enajenar ese derecho a terceros. \\

\textbf{¿Cómo puede usarse el derecho de superficie para financiar vivienda?} &
El titular del terreno constituye un derecho de superficie a favor de un desarrollador, quien construye unidades y las vende financiadas. Con lo cobrado recupera la inversión antes de que venza su propio plazo (30 a 50 años), lo que permite generar un derecho a la vivienda alternativo con seguridad jurídica. \\

\textbf{¿Sobre qué recae la hipoteca en ese esquema?} &
Sobre el propio derecho de superficie del comprador, nunca sobre el dominio del terreno, que permanece en cabeza de su titular original. Quien se asegura es el dueño del derecho de superficie. \\

\textbf{¿Qué ocurre si el derecho de superficie se subcontrata por un plazo menor?} &
Al extinguirse el plazo subcontratado, el derecho no vuelve al dueño del terreno sino al superficiario intermedio que lo había constituido, quien conserva el resto de su propio plazo. Nadie transmite más derechos de los que tiene. \\

\textbf{¿Cómo se extingue el derecho de superficie, en contraste con el usufructo?} &
El usufructo se extingue con la muerte del usufructuario (salvo derecho de acrecer pactado). El derecho de superficie se extingue por vencimiento del plazo o causales contractuales, no por fallecimiento del superficiario. \\

\textbf{¿Qué es el derecho de reversión y qué plazo tiene?} &
Cláusula de una donación que permite al donante recuperar el bien si el donatario fallece antes que él, siempre que ese fallecimiento ocurra dentro de diez años posteriores a la donación. Vencido ese plazo, la reversión no es exigible. \\

\textbf{¿Por qué la donación no siempre conviene como planificación sucesoria?} &
Puede generar iliquidez, exponer el bien a contingencias del donatario (por ejemplo, un divorcio) y dificultar el acceso al crédito del donatario, aun cuando el donante conserve el usufructo. Exige evaluación de alternativas. \\

\textbf{¿Qué riesgo caracteriza a la renta vitalicia?} &
Es un contrato aleatorio cuya conveniencia depende de un hecho incierto: cuánto vivirá el rentista. Puede ser un excelente o un pésimo negocio según variables que no pueden saberse de antemano. \\

\textbf{¿Para qué sirven las cláusulas de Drag-Along y Tag-Along?} &
El arrastre obliga a los minoritarios a vender junto con el mayoritario. El acompañamiento permite a los minoritarios sumarse a la venta en las mismas condiciones si lo desean. Protegen a ambas partes frente a cambios de control. \\

\textbf{¿Cómo se aplican esas cláusulas a un condominio?} &
Mediante un pacto entre condóminos que replique esa lógica, respaldado por un poder especial irrevocable. Los pactos solo son oponibles entre partes y a terceros que tomen conocimiento efectivo de ellos. \\

\textbf{¿Cómo se divide un condominio cuando alguien no quiere vender?} &
Si se trata de comunidad hereditaria, por vía de la división judicial de la herencia. Si se trata de un condominio de derecho real, por la acción de división de condominio, también judicial en caso de desacuerdo. \\

\textbf{¿Qué riesgo existe al comprar derechos sin consolidar?} &
Se puede estar adquiriendo apenas una expectativa o la cesión de un boleto sobre un proyecto sin planos aprobados, lo que puede derivar en modificaciones que perjudiquen al comprador. \\

\textbf{¿Desde cuándo se incorpora un heredero a una sociedad?} &
Hay tres posturas: desde el fallecimiento del causante, desde la declaratoria de herederos, o desde la partición. Es una cuestión difícil y en debate en el proyecto de reforma de la ley de sociedades. \\

\textbf{¿Qué diferencia hay entre quórum y mayoría?} &
El quórum es el porcentaje del capital con derecho a voto que debe estar presente para sesionar válidamente. La mayoría es el porcentaje de votos necesario, sobre los presentes, para aprobar una decisión. \\

\textbf{¿Qué diferencias fundamentales existen entre el derecho de uso y la habitación?} &
El uso es más amplio (incluye frutos, puede ser sobre muebles e inmuebles). La habitación es más limitada (solo personas humanas, solo inmuebles, derecho real únicamente). Ambos pueden ser vitalicios o temporales. \\

\textbf{¿Por qué preferir la habitación al usufructo para proteger a quien cuidó un familiar?} &
El usufructo permitiría alquilar sin habitar, desvirtuando la finalidad original. La habitación garantiza la morada efectiva y puede reconvertirse en sustento económico según circunstancias, sin perder su función de protección. \\

\midrule

\multicolumn{2}{p{\textwidth}}{

\textbf{SÍNTESIS FINAL}

La clase recorrió el conjunto de derechos reales que se desprenden del dominio para satisfacer necesidades de vivienda, financiamiento y planificación patrimonial sin recurrir necesariamente a la transmisión plena de la propiedad. El derecho de superficie aparece como una herramienta especialmente potente para el desarrollo inmobiliario: permite separar la titularidad del suelo de la titularidad de lo construido, constituir hipotecas sobre ese derecho y subcontratarlo por plazos menores. El usufructo, la reversión y la renta vitalicia plantean herramientas ligadas a planificación sucesoria y cuidado familiar, pero exigen análisis cuidadoso de riesgos. La donación puede generar iliquidez y complicaciones frente a terceros. Una idea transversal es que soluciones del derecho societario ---como arrastre y acompañamiento--- pueden trasladarse al condominio para evitar bloqueos entre copropietarios. La comunidad hereditaria, aunque funcione como condominio en la práctica, exige vías de división distintas. El derecho de uso y de habitación, aunque menos utilizados, resultan de gran utilidad práctica para garantizar vivienda de personas vulnerables, con solidez jurídica muy superior a acuerdos personales. En conjunto, la clase invita a una idea de fondo: antes de aplicar automáticamente figuras jurídicas tradicionales, conviene analizar la situación concreta de cada persona o familia y diseñar la combinación de derechos y cláusulas que mejor proteja sus intereses reales.

} \\

\bottomrule

\end{tabularx}
\end{table}

%=============================================================================
% FIN DEL DOCUMENTO
%=============================================================================

\end{document}
